\section{Applying the local theorems to the arithmetic family}
\label{sec:family-join}

\noindent\textbf{Status.}
The following arguments identify the actual canonical deficit, cyclic
coordinates and finite-place application. The eventual numerical suppliers
are proved in GI/GII, and the coefficient extension in BC. Existence of
the full global arithmetic construction retains the boundaries stated
below. The source's global chain length $m_{\rm ch}=D.R.m$ is distinct
from the local parameter $m_{\rm loc}=k-1$.
Here $\lceil z\rceil_+=\max(0,\lceil z\rceil)$ denotes the source's
natural-valued ceiling.

\paragraph{GJ1: the universal infimum is attained.}
Fix the orbit data and degree families. Let $S$ be the nonnegative
deficits for which the required certificate holds for every permitted
coefficient setting and dominant weight, at the source height
$\lceil t_{If}\rceil_+$, depth $D_{If}$ and conductor $8(D_{If}+1)$.
The zero certificate gives $1\in S$, so $q_0=\inf S\in[0,1]$.
Fix any one of those settings, weights and target lattices, and put
$T=(1-q_0)d$, where $d$ is its representation dimension.
For $T>0$ choose $q\in S$ sufficiently close to $q_0$ that
$(1-q)d>\lceil T\rceil-1$. Its integer certificate rank is at least
$\lceil T\rceil\ge T$. For $T\le0$ use rank zero. All other conditions
are unchanged. The fixed objects were arbitrary, so $q_0\in S$
with the original universal quantifiers. Thus the actual
$X.\texttt{kapI}=q_0$ supplies \texttt{adm7}, and is bounded above by
each admissible deficit, as required by GI's mean comparison.
This is the actual family choice: for $x$ beyond the family start,
the field \texttt{(Fm.ex x hx i).kapI\_def} says precisely that
\texttt{Fm.Gx.kapI} is this infimum. No replacement of the family
or its degree assignments is required.

\paragraph{GJ2: uniform scale estimates.}
Write $\varepsilon$ for the fixed family exponent and recall
\[
 a=x e^{x^\varepsilon},\quad X'=e^{x^{\varepsilon/3}},\quad
 b'=a/X',\quad \delta_1=x^{-\varepsilon/8},\quad
 p_0=k^2x^{26},\quad Y=x^{13}.
\]
GII2 gives $G'\le C_e\log x/\delta_1$ and
$1\le N_1\le C_e\log x$, uniformly in the family index.
Let $c_\Lambda=K.\texttt{La}>0$ and
$\Lambda=c_\Lambda N_1^2\log a$. For any $\eta>0$ choose $x\ge3$ with
\[
 \left(2c_\Lambda C_e^2+38C_e+c_\Lambda^{-1}\right)(\log x)^3
 \le x^\eta.
\]
This is an explicit uniform choice: $\log x\ge
\max(1,24M/\eta^4)$ suffices for the coefficient $M$ on the left,
by the fourth term of the exponential series. Since $k\le x^6$,
$\log a=\log x+x^\varepsilon$ and $p_0\le x^{38}$, we obtain
\begin{align*}
 x^{\varepsilon-\eta}&\le\Lambda\le x^{\varepsilon+\eta},\\
 p<p_0&\ \Longrightarrow\ \log(p^{G'})\le x^{\varepsilon/8+\eta},\\
 p<X'&\ \Longrightarrow\ \log(p^{G'})\le x^{11\varepsilon/24+\eta}.
\end{align*}
For the last two bounds use $\log p\le38\log x$ or
$\log p<x^{\varepsilon/3}$ respectively. These prove the actual
\texttt{adm5} fields; later we use $\eta=\varepsilon/48$.

\paragraph{GJ3: exact orbit identification.}
For a class-zero prime let $h=h_P$, $n'=\operatorname{ord}(h)$,
and let $t_0$ be its chosen orbit representative. The selector identities
$b_h(t+h)=b_h(t)$ and $t\in b_h(t)+\mathbb Nh$ imply
$b_h(b_h(t))=b_h(t)$. Hence
\[
 \phi(i)=t_0+ih\quad(0\le i<n')
\]
is a bijection from the cyclic positions onto this orbit. It identifies
$T=\{i:\phi(i)\in TT\}$ with $\{t\in TT:b_h(t)=t_0\}$, and therefore
\[
 N_v=|T|,\qquad a_{\rm loc}=a_L\circ\phi,\quad
 b_{\rm loc}=b_L\circ\phi,\qquad
 \sum_{t_0\in\mathrm{goodReps}}N_v\le|TT|. \tag{GJ.1}
\]
The last inequality is disjointness of the representative fibres.
The actual \texttt{GatTies}/\texttt{LatTied} equations identify these
positions and extremes with the coefficient representation at the place.
Its integer bound on the remaining coordinates is at most
$\lfloor c_G\rfloor$. Thus the local base-weight and dominance conditions
hold for precisely that representation.

\paragraph{GJ4: the two prime ranges give the exact orbit certificate.}
Assume \texttt{HypIAt}, \texttt{HypIIAt}, the local interfaces
\texttt{adm6}/\texttt{adm7}, $\rho\le\rho_B$, and the compatible
standard coefficient subring supplied by BC. Fix a served orbit at an
actual finite place. It is a good orbit of a candidate prime, so (GJ.1)
applies. Work first over the standard subring.

If $p<p_0$, apply the Type-II interface using \texttt{HypIIAt}; the
actual rank parameter is $r_G$. If $p\ge p_0$, the rank parameter is
one; \texttt{HypIAt} and GJ1 give the Type-I certificate. The candidate
condition excludes the band of width $\rho_B$, hence also the contained
band of width $\rho$. In both cases the definitions give exactly
\[
 t=\lceil t_v\rceil_+,\qquad D=D_p,\qquad
 \kappa=\kappa_v,\qquad c_*=8(D_p+1).                 \tag{GJ.2}
\]
BC transports the certificate to the actual coefficient DVR with
these same quantities. The image of the chosen square root of $p$
may have larger valuation than a uniformizer; the $p$-normalized
depth and the powers defining height and vanishing remain unchanged.
This proves the literal \texttt{OrbitLatticeBounds}. Empty served sets
require no certificate and supply no assertion of abundance.

\paragraph{GJ5: the uniform family consequence.}
The GI/GII bounds are eventual uniformly in the family index. Taking
the maximum of their finitely many thresholds and the family's start
therefore applies GJ4 to all members at once, whenever the actual
standard-subring and lattice identifications hold. One sufficient
Type-I coefficient vector is
\[
\begin{gathered}
 K.\texttt{La}\ge880,\quad K.\texttt{kI},K.\texttt{DI}\ge2800001,
 \quad K.\texttt{rho}\ge5600002,\\
 K.\texttt{kappaThresholdOne}\ge604228,\qquad
 K.\texttt{kappaThresholdTwo}\ge302114,\qquad C_m=302114.
\end{gathered}                                                   \tag{GJ.3}
\]
These constants precede all family members. The Type-II coefficient
conditions are recorded in GII. The source construction still has to
realize the selected subset, degree pair, local comparison data and the
remaining global objects together. GII proves a uniform Bernoulli
selection; the host also needs its biased coset selection and
\texttt{Frame.Shares}. GW proves the degree choice, and HW follows its
actual preselection and family transport. BC's ring identifications
come directly from supplied \texttt{GroupCompatibilityFamily} data;
\texttt{IsRecCEv} is an alternative sufficient interface. None of these
implications supplies the entire global construction by itself.
